% Einheit 1 von 4 – Seite berechnen mit Sinus, Kosinus und Tangens (bank/trigonometrie/e1.jsonl, zusammenbau v0.1)
%% TODO Zweigzeile Teil 1: was hier gelernt wird (Satz in Schülersprache aus dem Einheitstitel) – Einheit 1
\einheitenkopf[e1]{Einheit 1 von 4 · Seite berechnen mit Sinus, Kosinus und Tangens}
\zweigzeile{ab Kl. 10 (am Gymnasium ab Kl. 9) · P10 oft}
\setcounter{aufgabe}{11}

%% TODO Titel: Ich-kann-Satz fehlt in der Bank (Platzhalter: Kettenname) – Nr. 12
%% TODO Anweisung: ein Satz über der ersten Teilaufgabe fehlt in der Bank – Nr. 12
\begin{aufgabe}{sin, cos oder tan?}
\begin{teile}
\teil Rechter Winkel bei $B$. Vom Winkel $\alpha$ aus: $a$ ist gegeben, $b$ ist gesucht. Welche Funktion verbindet $a$ und $b$? \\ \kreuz{sin} \kreuz{cos} \kreuz{tan}
\end{teile}
\end{aufgabe}

%% TODO Titel: Ich-kann-Satz fehlt in der Bank (Platzhalter: Kettenname) – Nr. 13
%% TODO Anweisung: ein Satz über der ersten Teilaufgabe fehlt in der Bank – Nr. 13
\begin{aufgabe}{Mal oder geteilt?}
\begin{teile}
\teil $\mathrm{sin}\,23^\circ = \frac{x}{8}$, gesucht ist die Seite $x$. Wird mal genommen oder geteilt? \\ \kreuz{im Zähler: mal} \\ \kreuz{im Nenner: geteilt}
\end{teile}
\end{aufgabe}

%% TODO Titel: Ich-kann-Satz fehlt in der Bank (Platzhalter: Kettenname) – Nr. 14
%% TODO Anweisung: ein Satz über der ersten Teilaufgabe fehlt in der Bank – Nr. 14
\begin{aufgabe}{Seite berechnen}
\swz[3]{Rechter Winkel bei $B$. Schreibe an jede Seite H, G oder A – vom Winkel $\alpha$ aus.}{\dreieck{(0,0)}{(5,0)}{(5,3)}{r}{t}{s}{\alpha}{}{}}
\swz{Rechter Winkel bei $C$. Trage den Bruch ein. \hfill (P10 2017 OS)}{\dreieck{(0,0)}{(5,0)}{(1.8,2.4)}{d}{e}{f}{}{\beta}{}}
\swz{Rechter Winkel bei $B$. Trage den Bruch ein. \hfill (P10 2019 OS)}{\dreieck{(0,0)}{(5,0)}{(5,3)}{k}{m}{n}{\alpha}{}{}}
\swz{Rechter Winkel bei $A$. Trage den Bruch ein.}{\dreieck{(0,0)}{(4,0)}{(0,3)}{x}{y}{z}{}{}{\gamma}}
\swz{Rechter Winkel bei $C$. Trage den Bruch ein.}{\dreieck{(0,0)}{(5,0)}{(1.8,2.4)}{a}{b}{c}{\alpha}{}{}}
\swz{Rechter Winkel bei $A$. Trage den Bruch ein.}{\dreieck{(0,0)}{(4,0)}{(0,3)}{m}{k}{l}{}{\beta}{}}
\swfrage{Was bleibt gleich, was ändert sich?}
\begin{teile}
\teil Rechter Winkel bei $A$. Welche Gleichung gilt? \hfill (P10 2018 OS) \\ \kreuz{$\mathrm{sin}\,\beta = \frac{p}{q}$} \\ \kreuz{$\mathrm{cos}\,\beta = \frac{q}{p}$} \\ \kreuz{$\mathrm{sin}\,\beta = \frac{q}{p}$}
\end{teile}
%% TODO Darstellung für schwach fehlt in der Bank (trigonometrie-e1-k3-s3-v1; Abschnitt „Für schwache Schüler“)
\swz{Taschenrechner im Grad-Modus: $\mathrm{sin}\,23^\circ$ auf drei Dezimalen?}{}
%% TODO Darstellung für schwach fehlt in der Bank (trigonometrie-e1-k3-s4-v1; Abschnitt „Für schwache Schüler“)
\swz{Hypotenuse $8$ cm, $\alpha = 31^\circ$ – wie lang ist die Gegenkathete $x$ von $\alpha$?}{}
%% TODO Darstellung für schwach fehlt in der Bank (trigonometrie-e1-k3-s5-v1; Abschnitt „Für schwache Schüler“)
\swz{Hypotenuse $9$ cm, $\alpha = 27^\circ$ – wie lang ist die Ankathete $x$ von $\alpha$?}{}
\swz{Rechter Winkel bei $A$, $a = 7$ cm, $\beta = 41^\circ$. Wie lang ist $b$?}{\dreieck{(0,0)}{(4,0)}{(0,3)}{a}{b}{c}{}{\beta}{}}
%% TODO Darstellung für schwach fehlt in der Bank (trigonometrie-e1-k3-s7-v1; Abschnitt „Für schwache Schüler“)
\swz{Gegenkathete $6$ cm, $\alpha = 44^\circ$ – wie lang ist die Hypotenuse $x$?}{}
%% TODO Darstellung für schwach fehlt in der Bank (trigonometrie-e1-k3-s8-v1; Abschnitt „Für schwache Schüler“)
\swz{Ankathete $7$ cm, $\alpha = 32^\circ$ – wie lang ist die Hypotenuse $x$?}{}
%% TODO Darstellung für schwach fehlt in der Bank (trigonometrie-e1-k3-s9-v1; Abschnitt „Für schwache Schüler“)
\swz{Ankathete $6$ cm, $\alpha = 27^\circ$ – wie lang ist die Gegenkathete $x$?}{}
%% TODO Darstellung für schwach fehlt in der Bank (trigonometrie-e1-k3-s10-v1; Abschnitt „Für schwache Schüler“)
\swz{Gegenkathete $5$ cm, $\alpha = 31^\circ$ – wie lang ist die Ankathete $x$?}{}
%% TODO Darstellung für schwach fehlt in der Bank (trigonometrie-e1-k3-s11-v1; Abschnitt „Für schwache Schüler“)
\swz{Hypotenuse $9$ cm, $\alpha = 72^\circ$ – wie lang ist die Ankathete $x$? Wähle die Funktion selbst.}{}
%% TODO Darstellung für schwach fehlt in der Bank (trigonometrie-e1-k3-s12-v1; Abschnitt „Für schwache Schüler“)
\swz{$\text{$\mathrm{sin}\,48^\circ = \frac{6}{x}$ – wie groß ist $x$?}$}{}
%% TODO Darstellung für schwach fehlt in der Bank (trigonometrie-e1-k3-s13-v1; Abschnitt „Für schwache Schüler“)
\swz{Hypotenuse $6{,}8$ m, $\alpha = 42{,}6^\circ$ – wie lang ist die Gegenkathete $x$?}{}
%% TODO Darstellung für schwach fehlt in der Bank (trigonometrie-e1-k3-s14-v1; Abschnitt „Für schwache Schüler“)
\swz[3]{Hypotenuse $7$ cm, $\alpha = 42^\circ$. Zeige rechnerisch, dass die Ankathete $x$ rund $5{,}2$ cm lang ist.}{}
\swz{Eine $4{,}2$ m lange Leiter lehnt an einer Wand und bildet mit dem Boden einen Winkel von $71^\circ$ (Skizze). In welcher Höhe liegt sie an der Wand an?}{\dreieck{(0,0)}{(1.4,0)}{(1.4,4)}{Wand}{Leiter}{Boden}{71^\circ}{}{}}
\swz[3]{Ein Grundstück $ABCD$ hat bei $A$ einen rechten Winkel. Die Diagonale $BD$ ist $840$ m lang, der Winkel $ADB$ beträgt $47^\circ$. Zeige rechnerisch, dass $AB \approx 614$ m ist. \hfill (P10 2021 OS)}{}
\end{aufgabe}

%% TODO Titel: Ich-kann-Satz fehlt in der Bank (Platzhalter: Kettenname) – Nr. 15
%% TODO Anweisung: ein Satz über der ersten Teilaufgabe fehlt in der Bank – Nr. 15
\begin{aufgabe}{Seiten vom zweiten spitzen Winkel aus neu benennen (Gegen- und Ankathete tauschen)}
%% TODO Darstellung für schwach fehlt in der Bank (trigonometrie-e1-k4-s1-v1; Abschnitt „Für schwache Schüler“)
\swz{Rechter Winkel bei $C$. Vom Winkel $\alpha$ aus ist $a$ die Gegenkathete und $b$ die Ankathete. Wie heißen $a$ und $b$ vom Winkel $\beta$ aus?}{}
\end{aufgabe}

%% TODO Titel: Ich-kann-Satz fehlt in der Bank (Platzhalter: Kettenname) – Nr. 16
%% TODO Anweisung: ein Satz über der ersten Teilaufgabe fehlt in der Bank – Nr. 16
\begin{aufgabe}{Seitenverhältnis aus gegebenen Längen als Dezimalzahl und Vergleich mit dem Taschenrechnerwert des Winkels}
%% TODO Darstellung für schwach fehlt in der Bank (trigonometrie-e1-k5-s1-v1; Abschnitt „Für schwache Schüler“)
\swz{Gemessen: Gegenkathete $3{,}4$ cm, Hypotenuse $5$ cm, Winkel $43^\circ$. Verhältnis als Dezimalzahl – passt es zu $\mathrm{sin}\,43^\circ$?}{}
\end{aufgabe}

%% TODO Titel: Ich-kann-Satz fehlt in der Bank (Platzhalter: Kettenname) – Nr. 17
%% TODO Anweisung: ein Satz über der ersten Teilaufgabe fehlt in der Bank – Nr. 17
\begin{aufgabe}{Ergebnis prüfen: Kathete kürzer als Hypotenuse}
%% TODO Darstellung für schwach fehlt in der Bank (trigonometrie-e1-k6-s1-v1; Abschnitt „Für schwache Schüler“)
\swz[3]{Lea rechnet zur Hypotenuse $6$ cm eine Gegenkathete von $7{,}3$ cm aus. Kann das stimmen?}{}
\end{aufgabe}

%% TODO Titel: Ich-kann-Satz fehlt in der Bank (Platzhalter: Kettenname) – Nr. 18
%% TODO Anweisung: ein Satz über der ersten Teilaufgabe fehlt in der Bank – Nr. 18
\begin{aufgabe}{Seite berechnen – Fehler finden, Begründen, Anwendung}
\swz[3]{Gegenkathete $5$ cm, $\alpha = 41^\circ$, gesucht die Hypotenuse $x$. Wo ist der Fehler? \rechnung{\mathrm{sin}\,41^\circ &= \frac{5}{x} \\ x &= 5 \cdot \mathrm{sin}\,41^\circ \\ x &\approx 3{,}3 \text{ cm}}}{}
\swfrage{Warum ist $\mathrm{sin}\,\alpha$ im rechtwinkligen Dreieck immer kleiner als $1$?}
\swz[3]{Eine $6$ m lange Leiter soll mit $72^\circ$ zum Boden an der Hauswand stehen. Das Fenster beginnt in $5{,}5$ m Höhe. Reicht die Leiter bis zum Fenster?}{}
\end{aufgabe}

\uebersichtskasten{Seiten vom Winkel aus: Die Hypotenuse liegt dem rechten Winkel gegenüber, die Gegenkathete liegt dem markierten Winkel gegenüber, die Ankathete liegt an ihm – wechselt der Winkel, tauschen Gegen- und Ankathete. \\ Sinus, Kosinus, Tangens: sin α = Gegenkathete/Hypotenuse, cos α = Ankathete/Hypotenuse, tan α = Gegenkathete/Ankathete – der Taschenrechner (Grad-Modus) gibt den Wert zum Winkel. \\ Hypotenuse 12 cm, α = 40°: Gegenkathete = 12 · sin 40° ≈ 7,7 cm \quad Ankathete = 12 · cos 40° ≈ 9,2 cm \\ Steht die gesuchte Seite im Nenner, wird geteilt: Gegenkathete 9 cm, α = 40°: Hypotenuse = 9 : sin 40° ≈ 14,0 cm \quad Ankathete = 9 : tan 40° ≈ 10,7 cm \\ Merkregel: Kathete gesucht → Hypotenuse mal. Hypotenuse gesucht → Kathete geteilt. Erst am Ende runden.}
